\documentclass[11pt,usenames,dvipsnames,aspectratio=169]{beamer}

\usetheme[outer/progressbar=foot]{metropolis}
\usepackage{tikz}
\usepackage{graphicx}
\usepackage{xcolor}
\usepackage[utf8]{inputenc}
\usepackage[T1]{fontenc}
\usepackage[polish]{babel}
\usepackage{polski}
\usepackage{amsmath}
\usepackage{amssymb}

\setbeamercolor{palette primary}{bg=BlueViolet,fg=white}
\setbeamercolor{palette secondary}{bg=BlueViolet,fg=white}
\setbeamercolor{title separator}{bg=Red,fg=BlueViolet}
\setbeamercolor{progress bar}{bg=OrangeRed,fg=BlueViolet}

\usetikzlibrary{shapes.multipart}
\usetikzlibrary{decorations.pathreplacing}

\title{Dekompozycja szeregów czasowych}
\subtitle{ASC 2026}
\author{Piotr Żoch}
\date{}
\beamertemplatenavigationsymbolsempty

\begin{document}

\begin{frame}[plain]
\vspace*{-0.6cm}
\titlepage
\end{frame}

%--------------------------------------------------------------
\begin{frame}{Plan}
\begin{itemize}
    \item Składowe szeregów czasowych.
    \item Dekompozycja addytywna i multiplikatywna.
    \item Obserwacje nietypowe.
    \item Średnie ruchome.
    \item Metody wyrównywania sezonowego: X-13, TRAMO/SEATS, STL.
\end{itemize}
\end{frame}

%--------------------------------------------------------------
\section{Składowe szeregów czasowych}

\begin{frame}{Składowe szeregów czasowych}
\begin{center}
\includegraphics[scale=0.4]{pasted2.png}
\end{center}
\end{frame}

%--------------------------------------------------------------
\begin{frame}{Składowe szeregów czasowych}
\begin{center}
\includegraphics[scale=0.4]{pasted3.png}
\end{center}
\end{frame}

%--------------------------------------------------------------
\begin{frame}{Składowe szeregów czasowych}
\begin{itemize}
    \item Składowa przypadkowa.
    \item Składowa systematyczna (efekt oddziaływań stałego zestawu czynników na zmienną prognozowaną):
    \begin{itemize}
        \item Trend (tendencja rozwojowa).
        \item Stały poziom (przeciętny długookresowy poziom zmiennej).
        \item Składowa okresowa:
        \begin{itemize}
            \item Wahania sezonowe.
            \item Wahania cykliczne.
        \end{itemize}
    \end{itemize}
\end{itemize}
\end{frame}

%--------------------------------------------------------------
\begin{frame}{Składowe szeregów czasowych}
\begin{itemize}
    \item \textbf{Tendencja rozwojowa (trend)}
    \begin{itemize}
        \item Długookresowa skłonność do jednokierunkowych zmian (wzrostu lub spadku) wartości badanej zmiennej.
    \end{itemize}
    \item \textbf{Stały (przeciętny) poziom}
    \begin{itemize}
        \item Występuje, gdy w szeregu czasowym nie ma tendencji rozwojowej, wartości zmiennej prognozowanej oscylują zaś wokół pewnego (stałego poziomu).
    \end{itemize}
    \item \textbf{Wahania cykliczne}
    \begin{itemize}
        \item Długookresowe rytmiczne wahania wartości wokół tendencji rozwojowej lub stałego (przeciętnego) poziomu.
    \end{itemize}
    \item \textbf{Wahania sezonowe}
    \begin{itemize}
        \item Wahania wartości obserwowanej zmiennej wokół tendencji rozwojowej lub stałego (przeciętnego) poziomu tej zmiennej. Mają skłonność do powtarzania się w określonym czasie.
    \end{itemize}
\end{itemize}
\end{frame}

%--------------------------------------------------------------
\begin{frame}{Składowe szeregów czasowych}
\begin{center}
\includegraphics[scale=0.35]{pasted4.png}
\end{center}
\end{frame}

%--------------------------------------------------------------
\begin{frame}{Dekompozycja szeregów czasowych}
\begin{itemize}
    \item Wyodrębnienie poszczególnych składników nie zawsze jest łatwe z uwagi na rozmaite interakcje, które zachodzą pomiędzy nimi.
    \item Różne metody dekompozycji (średnie ruchome, wydzielanie komponentów o różnej częstotliwości, procedury TRAMO/SEATS itd).
\end{itemize}
\end{frame}

%--------------------------------------------------------------
\section{Trend i sezonowość}

\begin{frame}{Trend}
\begin{itemize}
    \item Trend
    \begin{itemize}
        \item \textbf{deterministyczny}
        \begin{align*}
        y_{t} & =\alpha+\beta t+\epsilon_{t}\\
        y_{t} & =\alpha+\beta t+\gamma t^{2}+\epsilon_{t}\\
        y_{t} & =\kappa e^{\delta t}e^{\epsilon_{t}}\\
        \epsilon_{t} & \sim iid\left(0,\sigma_{\epsilon}^{2}\right)
        \end{align*}
        \item \textbf{stochastyczny}
        \begin{align*}
        y_{t} & =\alpha+y_{t-1}+\epsilon_{t}\\
        \epsilon_{t} & \sim iid\left(0,\sigma_{\epsilon}^{2}\right)
        \end{align*}
    \end{itemize}
\end{itemize}
\end{frame}

%--------------------------------------------------------------
\begin{frame}{Trend}
\begin{center}
\includegraphics[scale=0.3]{pasted6.png}
\end{center}
\end{frame}

%--------------------------------------------------------------
\begin{frame}{Trend}
\begin{itemize}
    \item W praktyce rozróżnienie może nie być łatwe\ldots
\end{itemize}
\begin{center}
\includegraphics[scale=0.5]{pasted5.png}
\end{center}
\end{frame}

%--------------------------------------------------------------
\begin{frame}{Wahania sezonowe}
\begin{itemize}
    \item Dwa typy sezonowości:
    \begin{itemize}
        \item \textbf{Deterministyczna}
        \begin{itemize}
            \item Wartości obserwacji szeregu cechującego się sezonowością deterministyczną wahają się z amplitudą względnie stałą.
            \item Proces, którego bezwarunkowa średnia zależy od podokresu roku (np.\ miesiąca, kwartału).
            \item Ten rodzaj sezonowości jest modelowany za pomocą zmiennych zerojedynkowych, ich liczba zależna od liczby okresów w roku.
        \end{itemize}
        \item \textbf{Stochastyczna}
        \begin{itemize}
            \item Charakteryzuje się zmiennym w czasie wzorcem sezonowości.
        \end{itemize}
    \end{itemize}
\end{itemize}
\end{frame}

%--------------------------------------------------------------
\begin{frame}{Wahania sezonowe}
\begin{itemize}
    \item Deterministyczna sezonowość (przykład):
    \[
    y_{t}=\sum_{s=1}^{S}\delta_{t,s}\,m_{s}+\epsilon_{t}
    \]
    $\delta_{t,s}=1$ jeśli okres $t$ przypada na sezon $s$ (np.\ $\delta_{t,12}=1$ jeśli okres $t$ to grudzień w przypadku danych miesięcznych).
    \item Stochastyczna sezonowość (przykład -- sezonowy pierwiastek jednostkowy):
    \[
    y_{t}=y_{t-S}+\epsilon_{t}
    \]
\end{itemize}
\end{frame}

%--------------------------------------------------------------
\section{Dekompozycja addytywna i multiplikatywna}

\begin{frame}{Metody wyrównywania sezonowego}
\begin{itemize}
    \item Punktem wyjścia jest założenie, że szereg czasowy można zdekomponować na elementy:
    \begin{itemize}
        \item \textbf{czynnik sezonowy} $S_{t}$ -- regularne fluktuacje sezonowe;
        \item \textbf{cykl i trend} $C_{t}$ -- długookresowy trend, cykl koniunkturalny i inne składniki cykliczne;
        \item \textbf{trading-day} $TD_{t}$ -- kształt kalendarza -- korekta ze względu na różną długość miesięcy, liczbę dni roboczych itp.;
        \item \textbf{składnik losowy} $I_{t}$.
    \end{itemize}
    \item Chcemy uzyskać dane wyrównane sezonowo $SA_{t}$.
\end{itemize}
\end{frame}

%--------------------------------------------------------------
\begin{frame}{Postać addytywna}
\begin{itemize}
    \item \textbf{Postać addytywna}
    \begin{align*}
    Y_{t} & =S_{t}+C_{t}+TD_{t}+I_{t}\\
    SA_{t} & =C_{t}+I_{t}
    \end{align*}
    czyli amplituda wahań sezonowych i losowych jest stała.
    \item Czynnik sezonowy $S_t$ wyraża się w jednostkach zmiennej $Y_t$ (np.\ mld PLN).
    \item Warunek normalizacji: $\sum_{s=1}^{S}S_{t_s}=0$ w obrębie każdego roku.
\end{itemize}
\end{frame}

%--------------------------------------------------------------
\begin{frame}{Postać multiplikatywna}
\begin{itemize}
    \item \textbf{Postać multiplikatywna}
    \begin{align*}
    Y_{t} & =S_{t}\times C_{t}\times TD_{t}\times I_{t}\\
    SA_{t} & =C_{t}\times I_{t}
    \end{align*}
    czyli amplituda wahań sezonowych i losowych zależy od trendu.
    \item Czynnik sezonowy $S_t$ jest bezwymiarowy (np.\ $S_t=1{,}05$ oznacza, że wartość w tym sezonie jest o 5\% wyższa od trend-cyklu).
    \item Warunek normalizacji: $\sum_{s=1}^{S}S_{t_s}=S$ (średnia równa 1).
\end{itemize}
\end{frame}

%--------------------------------------------------------------
\begin{frame}{Postać addytywna vs postać multiplikatywna}
\begin{center}
\includegraphics[scale=0.5]{pasted7.png}
\end{center}
\end{frame}

%--------------------------------------------------------------
\begin{frame}{Postać addytywna vs postać multiplikatywna}
\begin{center}
\includegraphics[scale=0.4]{pasted8.png}
\end{center}
\end{frame}

%--------------------------------------------------------------
\begin{frame}{Transformacja logarytmiczna}
\begin{itemize}
    \item Model multiplikatywny można sprowadzić do addytywnego stosując logarytm:
    \begin{align*}
    Y_{t} & =S_{t}\times C_{t}\times I_{t}\\
    \ln Y_{t} & =\ln S_{t}+\ln C_{t}+\ln I_{t}
    \end{align*}
    \item Dlatego w praktyce często pracujemy na zlogarytmowanych danych.
    \item Dodatkowa zaleta: logarytm stabilizuje wariancję (jeśli wariancja rośnie proporcjonalnie do poziomu zmiennej).
    \item W kontekście wyrównywania sezonowego: wybór między dekompozycją addytywną a multiplikatywną jest jedną z pierwszych decyzji w procedurze.
\end{itemize}
\end{frame}

%--------------------------------------------------------------
\section{Obserwacje nietypowe}

\begin{frame}{Obserwacje nietypowe}
\begin{itemize}
    \item \textbf{O charakterze jednorazowym} \emph{(additive outlier (AO))} powodujące zmianę wartości zmiennej zależnej tylko w jednym okresie
    \[
    AO_{t}^{(t_{0})}=\begin{cases}
    1 & t=t_{0}\\
    0 & t\neq t_{0}
    \end{cases}
    \]
    \item Wpływ na szereg (addytywnie): $y_t = f(t) + \omega \cdot AO_t^{(t_0)} + \varepsilon_t$, gdzie $\omega$ to wielkość zaburzenia.
\end{itemize}
\end{frame}

%--------------------------------------------------------------
\begin{frame}{Obserwacje nietypowe}
\begin{itemize}
    \item \textbf{O charakterze przejściowym} \emph{(temporary change (TC))}, powodujące tymczasowe przesunięcie poziomu zmiennej zależnej, przy czym powrót ze skokowej zmiany wartości zmiennej zależnej do poziomu pierwotnego następuje zgodnie z funkcją wykładniczą w postaci
    \[
    TC_{t}^{(t_{0})}=\begin{cases}
    0 & t<t_{0}\\
    \alpha^{t-t_{0}} & t\geq t_{0}
    \end{cases}
    \]
    gdzie $\alpha\in(0,1)$ kontroluje szybkość powrotu (typowo $\alpha=0{,}7$).
\end{itemize}
\end{frame}

%--------------------------------------------------------------
\begin{frame}{Obserwacje nietypowe}
\begin{itemize}
    \item \textbf{O charakterze długotrwałym} \emph{(level shift (LS))}, powodujące trwałe przesunięcie poziomu zmiennej zależnej
    \[
    LS_{t}^{(t_{0})}=\begin{cases}
    0 & t<t_{0}\\
    1 & t\geq t_{0}
    \end{cases}
    \]
    \item W modelach regARIMA (TRAMO, X-13) obserwacje nietypowe są identyfikowane automatycznie i uwzględniane jako zmienne objaśniające:
    \[
    y_t = \sum_i \omega_i \cdot O_t^{(t_i)} + \underbrace{x_t'\beta + z_t}_{\text{regARIMA}} + \varepsilon_t
    \]
\end{itemize}
\end{frame}

%--------------------------------------------------------------
\section{Średnie ruchome}

\begin{frame}{Średnie ruchome}
\begin{itemize}
    \item Symetryczna średnia ruchoma rzędu $2q+1$:
    \[
    \hat{C}_{t}=\sum_{j=-q}^{q}w_{j}\,y_{t-j}
    \]
    gdzie $w_{j}=w_{-j}$ oraz $\sum_{j}w_{j}=1$.
    \item Aby wyeliminować sezonowość o okresie $S$, stosujemy średnią ruchomą o długości $S$ (np.\ 12 dla danych miesięcznych, 4 dla kwartalnych).
    \item Dla parzystego $S$ stosujemy średnią $2\times S$:
    \[
    \hat{C}_{t}=\frac{1}{S}\!\left(\frac{1}{2}y_{t-S/2}+y_{t-S/2+1}+\cdots+y_{t+S/2-1}+\frac{1}{2}y_{t+S/2}\right)
    \]
\end{itemize}
\end{frame}

%--------------------------------------------------------------
\begin{frame}{Średnie ruchome: właściwości}
\begin{itemize}
    \item Średnia ruchoma to \textbf{filtr liniowy}: przekształca szereg wejściowy $\{y_t\}$ na wyjściowy $\{\hat{C}_t\}$ za pomocą wag $\{w_j\}$.
    \item Średnia ruchoma o długości $S$ \textbf{eliminuje składową sezonową} o okresie $S$: jeśli $S_t$ jest periodyczna z okresem $S$ i $\sum_{s=1}^S S_{t_s} = 0$, to
    \[
    \frac{1}{S}\sum_{j=0}^{S-1} S_{t-j} = 0
    \]
    \item Jednocześnie zachowuje trend: jeśli $C_t = a + bt$, to $\hat{C}_t = a + bt$ (dla symetrycznych wag sumujących się do 1).
    \item Problem: średnia ruchoma nie daje wartości na końcach próby (brak obserwacji przyszłych) -- tu właśnie pomaga model ARIMA w X-13.
\end{itemize}
\end{frame}

%--------------------------------------------------------------
\begin{frame}{Filtr Hendersona}
\begin{itemize}
    \item W metodzie X-12/X-13 do estymacji trendu używa się \textbf{filtru Hendersona}, a nie prostej średniej ruchomej.
    \item Filtr Hendersona minimalizuje $\sum_{t}\!\left(\nabla^{3}\hat{C}_{t}\right)^{2}$, tj.\ sumę kwadratów trzeciej różnicy estymowanego trendu.
    \item Formalne sformułowanie: szukamy wag $\{w_j\}_{j=-q}^{q}$ rozwiązujących
    \[
    \min_{\{w_j\}} \sum_t \left(\Delta^3 \hat{C}_t\right)^2 \quad \text{p.w.} \quad \sum_j w_j = 1,\;\; \sum_j j^k w_j = 0 \;\text{ dla } k=1,2
    \]
    \item Typowe długości filtra: 9 lub 13 (dane miesięczne), 5 lub 7 (dane kwartalne).
    \item Wybór długości zależy od stosunku zmienności komponentu nieregularnego do zmienności trend-cyklu (\emph{I/C ratio}).
\end{itemize}
\end{frame}

%--------------------------------------------------------------
\section{Metody wyrównywania sezonowego}

\begin{frame}{Metody wyrównywania sezonowego}
\begin{itemize}
    \item Dwie popularne metody to:
    \begin{itemize}
        \item \textbf{X-12-ARIMA} (obecnie X-13-ARIMA-SEATS)
        \item \textbf{TRAMO/SEATS}
    \end{itemize}
    \item Obie dostępne w programie \textbf{JDemetra+} (opracowany przez Eurostat, darmowy).
\end{itemize}
\end{frame}

%--------------------------------------------------------------
\begin{frame}{X-13-ARIMA-SEATS: struktura}
\begin{itemize}
    \item Program składa się z dwóch etapów:
    \begin{enumerate}
        \item \textbf{regARIMA} -- wstępne oczyszczenie szeregu (patrz następny slajd).
        \item \textbf{X-11} lub \textbf{SEATS} -- właściwa dekompozycja na komponenty.
    \end{enumerate}
    \item Kolejność: najpierw regARIMA usuwa z szeregu efekty kalendarzowe i obserwacje nietypowe, a następnie X-11 lub SEATS dekomponuje oczyszczony szereg.
\end{itemize}
\end{frame}

%--------------------------------------------------------------
\begin{frame}{regARIMA}
\begin{itemize}
    \item Model regARIMA to model regresji, w którym reszty mają strukturę SARIMA:
    \[
    y_t = \underbrace{x_t'\beta}_{\text{regresory}} + \underbrace{z_t}_{\text{reszty SARIMA}}
    \]
    \item Regresory $x_t$ mogą obejmować: zmienne sezonowe, efekty kalendarzowe (trading-day, Wielkanoc), obserwacje nietypowe (AO, TC, LS).
    \item Reszty $z_t$ spełniają model SARIMA:
    \[
    \phi(B)\,\Phi(B^S)\,(1-B)^d(1-B^S)^D\, z_t = \theta(B)\,\Theta(B^S)\,\varepsilon_t
    \]
    \item Po estymacji modelu obliczamy \textbf{zlinearyzowany szereg}: $\tilde{y}_t = y_t - x_t'\hat{\beta}$, który trafia do X-11 lub SEATS.
    \item Model SARIMA służy też do \textbf{prognozowania} wartości na końcach próby (potrzebnych do filtrów symetrycznych).
\end{itemize}
\end{frame}

%--------------------------------------------------------------
\begin{frame}{X-11: algorytm iteracyjny}
\small
\begin{itemize}
    \item Metoda ad hoc, dobór filtrów nie zależy od właściwości statystycznych szeregu.
    \item Nieco myląca nazwa -- model ARIMA jest wykorzystywany jedynie do estymacji wartości prognozowanych.
    \begin{enumerate}
        \item Trend obliczany za pomocą średnich ruchomych (filtr $2\times 12$ lub $2\times 4$).
        \item Odjąć trend od szeregu, pozostaje komponent sezonowy i nieregularny: $\hat{S}_t + \hat{I}_t = Y_t - \hat{C}_t$.
        \item Komponent sezonowy obliczany za pomocą średnich ruchomych (sezonowe $3\times 3$, $3\times 5$ lub $3\times 9$).
        \item Odjąć komponent sezonowy od szeregu: $\hat{SA}_t = Y_t - \hat{S}_t$.
        \item Ponowić (1) i (2) z filtrem Hendersona dla trendu.
        \item Ponowić (3) i (4) z nieco innymi wagami.
        \item Otrzymać trend i komponent nieregularny odejmując uzyskany w~(6) komponent sezonowy od szeregu.
    \end{enumerate}
\end{itemize}
\end{frame}

%--------------------------------------------------------------
\begin{frame}{X-11: filtry sezonowe}
\begin{itemize}
    \item W kroku (3) stosujemy \textbf{sezonowe średnie ruchome}: osobna średnia ruchoma dla każdego miesiąca/kwartału.
    \item Np.\ filtr $3\times 3$ dla danych miesięcznych: dla każdego miesiąca $s$ stosujemy dwukrotnie prostą średnią z~trzech kolejnych lat, co daje ważoną średnią z~pięciu lat:
    \[
    \hat{S}_{t}^{(s)} = \frac{1}{9}\!\left(y_{t-24}^{SI}+2y_{t-12}^{SI}+3y_t^{SI}+2y_{t+12}^{SI}+y_{t+24}^{SI}\right)
    \]
    gdzie $y_t^{SI} = Y_t - \hat{C}_t$.
    \item Następnie normalizacja: $\hat{S}_t \leftarrow \hat{S}_t - \overline{\hat{S}}$ (addytywnie) lub $\hat{S}_t \leftarrow \hat{S}_t / \overline{\hat{S}}$ (multiplikatywnie).
    \item Wybór filtra ($3\times 3$, $3\times 5$, $3\times 9$) zależy od stabilności wzorca sezonowego: im bardziej zmienny, tym krótszy filtr.
\end{itemize}
\end{frame}

%--------------------------------------------------------------
\begin{frame}{TRAMO/SEATS}
\begin{itemize}
    \item \textbf{TRAMO} (``Time Series Regression with ARIMA Noise, Missing Observations and Outliers''): dopasowuje proces SARIMAX do szeregu i oczyszcza szereg z wpływu wybranych czynników $X$ (kalendarzowych, nietypowych\ldots) za pomocą tego modelu.
    \item \textbf{SEATS} (``Seasonal Extraction in ARIMA Time Series''): korzysta ze specyfikacji SARIMA, by zdekomponować szereg na elementy: trendu, sezonowości, nieregularny.
\end{itemize}
\end{frame}

%--------------------------------------------------------------
\begin{frame}{SEATS: dekompozycja oparta na modelu}
\begin{itemize}
    \item Załóżmy, że zlinearyzowany szereg ma model SARIMA:
    \[
    \phi(B)\,\Phi(B^S)\,(1-B)^d(1-B^S)^D\,y_t = \theta(B)\,\Theta(B^S)\,\varepsilon_t
    \]
    \item SEATS dekomponuje to na \textbf{nieobserwowalne komponenty}:
    \[
    y_t = C_t + S_t + I_t
    \]
    gdzie każdy komponent ma własny model ARIMA, np.:
    \begin{align*}
    \phi_C(B)\,(1-B)^d\,C_t &= \theta_C(B)\,\varepsilon_t^C \\
    \phi_S(B)\,(1-B^S)^D\,S_t &= \theta_S(B)\,\varepsilon_t^S
    \end{align*}
\end{itemize}
\end{frame}

%--------------------------------------------------------------
\begin{frame}{SEATS: ekstrakcja sygnału}
\begin{itemize}
    \item Ekstrakcja sygnału za pomocą \textbf{filtrów Wienera-Kołmogorowa} -- optymalnych filtrów liniowych minimalizujących błąd średniokwadratowy.
    \item Intuicja: mając model ARIMA całego szeregu i modele poszczególnych komponentów, możemy obliczyć optymalną estymację każdego komponentu warunkując na obserwacjach $\{y_1,\ldots,y_T\}$:
    \[
    \hat{C}_t = E\!\left[C_t \mid y_1, \ldots, y_T\right]
    \]
    \item Filtry te są \textbf{wyznaczone przez model ARIMA} -- nie wymagają arbitralnych wyborów (jak długość filtra w X-11).
    \item Kluczowa różnica wobec X-11: filtry dostosowane do właściwości statystycznych konkretnego szeregu (a nie ad hoc).
\end{itemize}
\end{frame}

%--------------------------------------------------------------
\begin{frame}{X-11 vs SEATS: porównanie}
\begin{itemize}
    \item \textbf{X-11}:
    \begin{itemize}
        \item Metoda nieparametryczna (filtry ad hoc).
        \item Nie wymaga specyfikacji modelu.
        \item Sprawdzona, stabilna procedura (używana od lat 60.).
    \end{itemize}
    \item \textbf{TRAMO/SEATS}:
    \begin{itemize}
        \item Metoda parametryczna (oparta na modelu SARIMA).
        \item Filtry optymalne w sensie MSE.
        \item Dostarcza błędy standardowe komponentów.
    \end{itemize}
    \item W praktyce wyniki obu metod są zazwyczaj bardzo zbliżone.
    \item X-13-ARIMA-SEATS łączy obie metody w jednym programie.
\end{itemize}
\end{frame}

%--------------------------------------------------------------
\begin{frame}{Dekompozycja STL}
\begin{itemize}
    \item \textbf{STL} (\emph{Seasonal and Trend decomposition using Loess}) -- Cleveland et al.\ (1990).
    \item Iteracyjna procedura wykorzystująca lokalną regresję wielomianową (loess) do estymacji trendu i sezonowości.
    \item Zalety:
    \begin{itemize}
        \item Elastyczność -- użytkownik kontroluje gładkość trendu i sezonowości.
        \item Odporność na obserwacje nietypowe (wersja robust).
        \item Może obsługiwać dowolną częstotliwość danych.
    \end{itemize}
    \item Ograniczenia:
    \begin{itemize}
        \item Tylko dekompozycja addytywna (multiplikatywną uzyskujemy przez logarytm).
        \item Nie ma wbudowanego modelowania efektów kalendarzowych ani obserwacji nietypowych (w odróżnieniu od X-13 i TRAMO/SEATS).
    \end{itemize}
\end{itemize}
\end{frame}

\end{document}
